\chapter{Verification and reproducibility}
\label{app:repro}

This appendix states what the executable checks of the manuscript compute, by which algorithms, under which conventions and with which
limitations, and maps every labelled statement to the checks that test it. The mathematics of the manuscript does not depend on this
appendix: every definition, theorem and proof is in the chapters or in \crefrange{app:A}{app:branching}, and a check is evidence for a
coded assertion, not a proof of the statement it tests. The scripts are supplementary material in the project repository, under \path{verification/chNN_<topic>/} (one directory per chapter, shared code in \path{verification/common/}) of the repository \url{https://github.com/driss-bencheikh/spacetime-electromagnetic-framework}, at the commit of version 1.0.0 recorded in \texttt{release/README.md}; this appendix gives what is needed to rerun and to judge them.

\section{Status vocabulary}
\label{app:repro:status}

A statement is \emph{analytically established} when its proof is in the text or it is cited, with its hypotheses, from published
literature (\emph{inherited}); \emph{exactly verified} when a check establishes the coded assertion in exact arithmetic (integers,
rationals, algebraic numbers, or symbolic expressions simplified to zero); \emph{numerically supported} when the check uses floating-point
arithmetic, a random sample or a numerical integration, with a stated tolerance; \emph{conditional} when it rests on an explicit
assumption; \emph{obstructed} or \emph{refuted} when a no-go or a counterexample is proved; \emph{unresolved} otherwise. A \texttt{PASS} line
certifies only the assertion printed with it: a check on one example verifies that example, and a check of an identity at a point does not
verify the theory built on it.

\section{Realisations and conventions}
\label{app:repro:conv}

All checks use the integer realisation of \cref{app:realisation}: $e_1=\sigma_x\otimes\id\otimes\id$, $e_2=\sigma_z\otimes\sigma_x\otimes\id$,
$e_3=\sigma_z\otimes\sigma_z\otimes\sigma_x$, $e_4=\varepsilon\otimes\id\otimes\id$, $e_5=\sigma_z\otimes\varepsilon\otimes\id$, $e_6=\sigma_z\otimes\sigma_z\otimes\varepsilon$
with $\varepsilon=\bigl(\begin{smallmatrix}0&-1\\1&0\end{smallmatrix}\bigr)$, so $e_a^2=\eta_{aa}\id$ with $\eta=\operatorname{diag}(1,1,1,-1,-1,-1)$, every
blade is a signed permutation matrix and $e_a^{\!\top}=\eta_{aa}e_a$. At $n$ constituents the generators are
$E_{6k+a}=\omega^{\otimes k}\otimes e_a\otimes\id^{\otimes(n-1-k)}$ with $\omega=e_1\cdots e_6$ ($64\times64$ at $n=2$, $512\times512$ at $n=3$). The
interior gamma matrices are $\gamma^\mu=(e_1,e_4,e_5,e_6)$ (or the clock four-plane $\langle\widehat W,e_4,e_5,e_6\rangle$ where a statement
says so), $\gamma_\mu=\eta_{\mu\nu}\gamma^\nu$ with $\eta=\operatorname{diag}(1,-1,-1,-1)$, $\Omega=\gamma^0\gamma^1\gamma^2\gamma^3$ and
$I_4=\gamma_0\gamma_1\gamma_2\gamma_3=-\Omega$; $J=e_2e_3$ at one constituent and the product of the internal generators at $n$; $\Gamma=J\Omega$;
two-forms are identified with bivectors by $F\mapsto\frac12F_{ab}\gamma^a\gamma^b$ and $\varepsilon_{0123}=+1$. The kinetic kernels are
$K^\mu=-\gamma^0\gamma^\mu$ at one constituent (\cref{prop:kernels1}) and $P_*\gamma^0\gamma^\mu P_*$ on the chiral block, an overall sign
of the fermion Lagrangian that changes no field equation, gauge condition or spectrum (\cref{r1:prop:dirac}).

\section{Algorithms}
\label{app:repro:alg}

\begin{enumerate}[label=(A\arabic*),leftmargin=2.6em]
\item \emph{Blade parity.} $e_Ie_J=(-1)^{|I||J|-|I\cap J|}e_Je_I$ (\cref{lem:blade}) decides commutation of blades without matrices; counts over
all $2^{18}$ blades of $\Cl(9,9)$ (\cref{thm:fourclasses}) are exact enumerations.
\item \emph{Exact linear solution spaces.} Commutants, intertwiners, invariant forms and gauge conditions ($X^{\!\top}K+KX=0$, $[X,L]=0$,
$aN+Na^{\!\top}=0$) are linear in the unknown matrix; the conditions are written on the blade basis (or on all matrix entries) and solved by
exact rational nullspaces. No ansatz is made: the unknown ranges over the whole matrix algebra.
\item \emph{Gram certificates on the large modules.} For operators $A_k$ on $\R^{512}$ the integer Gram matrix $\Tr(A_k^{\!\top}A_l)$ has the rank of
their span; commuting subspaces (the nullspace of $X\mapsto[X,P]$), kernels on a block and images are obtained from such Gram matrices in exact
arithmetic (\cref{thm:table}).
\item \emph{Modular ranks with explicit solutions.} Where a dimension is asserted for a large integer system, the rank modulo a large prime
bounds the rational rank from below; explicit exact solutions bound it from above, and the two are made to meet (e.g.\ \cref{thm:survive}).
Where only ranks modulo two primes are computed (\cref{thm:inertia3}), the check says so.
\item \emph{Symbolic geometry.} Curvature of the stated metrics by Christoffel symbols, Riemann, Ricci and Einstein tensors in exact
symbolic arithmetic, with closed-form inverses for Kaluza--Klein coframes; Euler--Lagrange expressions and Poisson brackets by symbolic
differentiation with general functions.
\item \emph{Spectra of exact integer matrices.} Characteristic polynomials are computed exactly where possible (\cref{thm:spectra}); joint
spectra of commuting integer matrices on the $512$-dimensional module are read by floating-point diagonalisation of a generic linear
combination and accepted only within $10^{-8}$ of the stated rationals (numerical).
\item \emph{Vacuum minimisation.} The radiative energy $E_x$ on the $272$-dimensional kernel space is minimised from random starts and its
critical points are classified by the Hessian (\cref{ch:ch07_vacuum}); the stated minimum is numerically supported and its global
minimality is a conjecture.
\item \emph{Integrations.} Ordinary differential equations (closure dynamics, orbits) are integrated with stated tolerances
(\texttt{DOP853} or \texttt{LSODA}, relative tolerance $10^{-12}$) and compared with closed forms where these exist.
\end{enumerate}

\section{Exact certificates}
\label{app:repro:cert}

The following identities are certified in exact arithmetic and carry the main structural statements: the Clifford relations of the
realisation at $n=1,2,3$; the trace-orthogonality $\Tr(\theta_S^{\!\top}\theta_R)=16\,\delta_{SR}$ of the $256$ monomials of the module of
\cref{thm:module}; the characteristic polynomials of \cref{thm:spectra}; the solution spaces of \cref{thm:C,thm:PT,prop:gauge1,prop:massangle}
over all of $M_8(\R)$; the commutants and invariant forms of \cref{cor:Gselected,prop:lGhoms,thm:inertiareq}; the grade census of
\cref{prop:censusint}; $\Tr P_*=64$ and the seven traceless products of \cref{thm:PS}; the projector-family table of \cref{thm:table}; the
$8192=4\times2048$ blades of \cref{thm:fourclasses}; the centraliser dimensions $72$, $30$ and the Casimir $3\,\id$ of
\cref{thm:so9C,thm:nocommon}; the curvature identities of \cref{ch:ch11_curvature} for general operators and the stated metrics; the
first-class bracket algebra of \cref{prop:constraints} for a general metric and potential. The release audit adds: the kernel spaces and
internal algebras of \cref{thm:subkernels} on the blades of $\Cl(3,3)$; the $CP$, $C$ and $P$ images of $\Phi_0$ in \cref{thm:vacCP}, with $\Phi_0$
rebuilt in integer arithmetic; the counts of \cref{thm:Bcounts} graded by $F$ and by the $J\Gamma$ charge; the centre and generators of the Lorentz
commutant (\cref{prop:config}); the basis of $\Pfam\mathcal C\Pfam$ by blade classes and the exact rank of the curvature values in
\cref{thm:holonomy,prop:chiralpreserved}; the second-rung family of \cref{thm:towervac,thm:S4,thm:towercurved,prop:towermatter}; and the
weights of the Weyl block at $n=3$ in \cref{prop:spectator} (a6, a7).

\section{Numerical checks and limitations}
\label{app:repro:num}

About one check in eight is numerical. They concern spectra on the large modules read in floating point (A6), random samples (random
clocks, random spinors, random points of projector families), numerical integrations (A8), the vacuum minimisation (A7), loop-level
coefficients, and the electroweak recomputations a8 on the kernel space, whose ranks are read with an asserted spectral gap. Each states its tolerance in its claim text. Their limitations are: a sample does not cover a family; a tolerance does not
exclude a nearby degenerate case; the minimisation finds the lowest critical points reached from the starts used and does not prove
global minimality; numerical integrations verify conserved quantities over finite time intervals only. Statements resting on such checks are
classed as numerically supported, not as exactly verified.

\section{Software and procedure}
\label{app:repro:soft}

Python~3 with \texttt{numpy}, \texttt{scipy}, \texttt{sympy} and \texttt{mpmath}; no other dependency. Each script prints one
\texttt{PASS} or \texttt{FAIL} line per check and exits with status zero exactly when all its checks pass. From the repository root:
{\raggedright
\begin{itemize}[leftmargin=1.6em]
\item \texttt{python3} \path{verification/run_all.py} \texttt{-{}-out} \path{<directory>} \texttt{-j 3} runs every script and writes one output file
per script, a summary and the software environment;
\item \texttt{python3} \path{verification/tools/make_manifest.py} \texttt{-{}-runs} \path{<directory>} rebuilds \path{verification/MANIFEST.tsv} from
the outputs; \path{verification/tools/coverage.py} maps statements to cited checks; \path{verification/tools/make_appendix_d.py} regenerates
the theorem-to-code table \path{verification/THEOREM_CODE.tsv} and the list of scripts \path{verification/SCRIPTS.tsv}; and \path{verification/tools/make_reproducibility_index.py} writes \path{reproducibility/manifest.json} and
the indexes of exact and numerical checks in \path{reproducibility/selected-results/};
\item the manuscript is built from \path{manuscript/} with \texttt{pdflatex}, \texttt{bibtex} and two further passes of \texttt{pdflatex} on
\path{main.tex}.
\end{itemize}\par}
The run records of this version are in \path{reproducibility/selected-results/v1.0.1/}; those of version 1.0.0 are kept in
\path{reproducibility/selected-results/v1.0.0/}.

The theorem-to-code table is kept in the repository, not printed: \path{verification/THEOREM_CODE.tsv} lists every labelled statement with
its environment, the checks that test it (cited in its text, added by the audit, or related without citation) and its evidence, which is
\emph{exact}, \emph{numerical} or \emph{exact+numerical} when checks exist and otherwise \emph{proof in text}, \emph{inherited}, \emph{conjecture}
or \emph{argument in text} (argued in the surrounding text without a proof environment or a check). \path{verification/SCRIPTS.tsv} lists the
scripts with their numbers of exact and numerical checks.
